% 江苏海洋大学毕业设计（论文）选题情况汇总表
% Word XML 对齐: Table 26 (15列，横向页面)

\documentclass[a4paper,landscape,zihao=-4]{ctexart}

\RequirePackage{geometry}
\RequirePackage{../../styles/jouhandbook}

\geometry{
    a4paper,
    landscape,
    top=2.2cm,
    bottom=2.2cm,
    left=2.3cm,
    right=2.3cm
}

\begin{document}
\pagestyle{empty}

{\zihao{-4}\songti
\begin{center}
    {\heiti 江苏海洋大学\hspace{2em}届毕业设计（论文）选题情况汇总表}
\end{center}

\vspace{0.2cm}

学院：\hspace{8cm}专业：

\vspace{0.2cm}

\begin{center}
\tiny
\renewcommand{\arraystretch}{1.1}
\begin{tabular}{|C{0.81cm}|C{5.74cm}|C{0.70cm}|C{0.70cm}|C{0.70cm}|C{0.70cm}|C{1.03cm}|C{1.91cm}|C{1.91cm}|C{1.91cm}|C{1.91cm}|C{1.81cm}|C{1.16cm}|C{1.81cm}|C{1.24cm}|}
\hline
\multicolumn{15}{|C{24.07cm}|}{{\heiti 江苏海洋大学\hspace{2em}届毕业设计（论文）选题情况汇总表}} \\
\hline
\multicolumn{15}{|L{24.07cm}|}{学院：\hspace{8cm}专业：} \\
\hline
\multirow{2}{*}{序号} & \multicolumn{6}{C{9.54cm}|}{毕业设计（论文）课题} & \multicolumn{3}{C{5.73cm}|}{指 导 教 师} & \multicolumn{2}{C{3.72cm}|}{学\hspace{1em}生} & \multirow{2}{*}{成绩} & \multirow{2}{*}{设计/论文} & \multirow{2}{*}{班级} \\
\cline{2-12}
 & 名\hspace{2em}称 & 性质 & 来源 & 难度 & 份量 & 结合实际 & 姓\hspace{1em}名 & 职\hspace{1em}称 & 工作单位 & 姓\hspace{1em}名 & 学\hspace{1em}号 &  &  &  \\
\hline
 & & & & & & & & & & & & & & \\
\hline
 & & & & & & & & & & & & & & \\
\hline
 & & & & & & & & & & & & & & \\
\hline
 & & & & & & & & & & & & & & \\
\hline
 & & & & & & & & & & & & & & \\
\hline
 & & & & & & & & & & & & & & \\
\hline
 & & & & & & & & & & & & & & \\
\hline
 & & & & & & & & & & & & & & \\
\hline
 & & & & & & & & & & & & & & \\
\hline
 & & & & & & & & & & & & & & \\
\hline
\multicolumn{15}{|L{24.07cm}|}{%
（请在最后一页按专业统计如下项目：）\par
总课题数：\hspace{1.5cm}（其中设计：\hspace{1cm}论文：\hspace{1cm}）；结合实际课题所占比例：\hspace{1.5cm}。\par
工程实践类课题占比例：\hspace{1cm}；实验研究课题占比例：\hspace{1cm}；理论研究课题占比例：\hspace{1cm}；其他课题占比例：\hspace{1cm}。%
} \\
\hline
\end{tabular}
\end{center}

\vspace{0.3cm}
\zihao{5}
注：1．“课题性质”栏用英文字母表达，分别为：A.工程实践，B.实验研究，C.理论研究，D.其他\par
2．“课题来源”栏用英文字母表达，分别为：A.指导教师的科研课题；B.指导教师收集的科研和生产实际中的课题；C.学生在科学活动和工程实践中自立的课题；D.自拟课题。\par
3．“难度”栏填写预计难度：A.难；B.一般；C.易。4．“份量”栏填写课题预计工作量：A.大；B.适中；C.小。\par
5．结合实际的课题在“结合实际”栏中打“√”。6．“设计/论文”栏填写 A.设计 B.论文。7．本表一式两份，一份由学院存档，一份交教务处实践科。\par

\vspace{0.5cm}
填表人：\hspace{6cm}学院负责人：\hspace{4cm}（公章）\hspace{1cm}年\hspace{0.6cm}月\hspace{0.6cm}日
}

\end{document}
